\documentclass[11pt,a4paper]{article}
\usepackage[T1]{fontenc}
\usepackage{lmodern,microtype}
\usepackage[margin=26mm,headheight=15pt]{geometry}
\usepackage{amsmath,amssymb,amsthm,mathtools}
\usepackage{booktabs,array,longtable,enumitem}
\usepackage[dvipsnames]{xcolor}
\usepackage{fancyhdr}
\usepackage{hyperref}
\hypersetup{colorlinks=true,linkcolor=MidnightBlue,citecolor=MidnightBlue,
 urlcolor=MidnightBlue,pdftitle={Finite-Jet Pressure Laws for Gevrey Reversion and Condensation},
 pdfauthor={ChatGPT; prepared for Vladimir Reshetnikov}}
\usepackage[nameinlink,noabbrev]{cleveref}
\crefname{lemma}{lemma}{lemmas}\Crefname{lemma}{Lemma}{Lemmas}
\crefname{proposition}{proposition}{propositions}\Crefname{proposition}{Proposition}{Propositions}
\crefname{corollary}{corollary}{corollaries}\Crefname{corollary}{Corollary}{Corollaries}
\crefname{definition}{definition}{definitions}\Crefname{definition}{Definition}{Definitions}
\crefname{example}{example}{examples}\Crefname{example}{Example}{Examples}
\setlength{\parindent}{0pt}
\setlength{\parskip}{4pt}
\setlength{\emergencystretch}{3em}
\pagestyle{fancy}\fancyhf{}
\fancyhead[L]{\small Finite-jet pressure laws}
\fancyhead[R]{\small ProveIt research continuation}
\fancyfoot[C]{\thepage}
\newtheorem{theorem}{Theorem}[section]
\newtheorem{lemma}[theorem]{Lemma}
\newtheorem{proposition}[theorem]{Proposition}
\newtheorem{corollary}[theorem]{Corollary}
\theoremstyle{definition}
\newtheorem{definition}[theorem]{Definition}
\newtheorem{example}[theorem]{Example}
\theoremstyle{remark}
\newtheorem{remark}[theorem]{Remark}
\newcommand{\N}{\mathbb N}
\newcommand{\R}{\mathbb R}
\newcommand{\C}{\mathbb C}
\newcommand{\E}{\mathbb E}
\newcommand{\PP}{\mathbb P}
\newcommand{\ee}{\mathrm e}
\newcommand{\inv}{^{\langle-1\rangle}}
\newcommand{\TV}{d_{\mathrm{TV}}}
\newcommand{\ind}{\mathbf 1}
\newcommand{\coef}[2]{[#1]\,#2}
\newcommand{\fall}[2]{#1^{\underline{#2}}}
\newcommand{\rise}[2]{(#1)_{#2}}
\DeclareMathOperator{\Pois}{Pois}
\DeclareMathOperator{\Var}{Var}
\DeclareMathOperator{\Law}{Law}
\DeclareMathOperator{\Rea}{Re}
\numberwithin{equation}{section}
\newcommand{\snapshot}{\texttt{a4268e78\,ebd0f6bf\,71ba609c\,07f4b341\,5748f532}}
\begin{document}
\hypersetup{pageanchor=false}
\begin{titlepage}
\vspace*{10mm}
{\color{MidnightBlue}\Large ProveIt research continuation\par}
\vspace{11mm}
{\Huge\bfseries Finite-Jet Pressure Laws\par}
\vspace{3mm}
{\huge\bfseries for Gevrey Reversion\par and Condensation\par}
\vspace{9mm}
{\Large Complete multiplicative asymptotics,\par explicit critical constants, and Poisson centering\par}
\vspace{11mm}
{\large Prepared for Vladimir Reshetnikov\par}
{\large ChatGPT\par}
{30 September 2026\par}
\vspace{9mm}
\begin{abstract}
The ProveIt research corpus leaves a precise gap in its sharp
fixed-coefficient-ball analysis of formal reversion: at Gevrey orders
$0<s<1/2$, the leading logarithmic amplification is known, but a
multiplicative equivalent is not supplied. We close that gap for the
stated factorial weights and, more generally, for positive weights with
an eventually shifted-factorial tail and an arbitrary finite head.
The logarithm of the normalized extremal inverse coefficient is a finite
sum $\sum_{ks\le1}p_k n^{1-ks}+o(1)$, with an explicit coefficient
extraction formula for every $p_k$. Thus every positive order, including
all reciprocal-integer transition points, has a complete multiplicative
equivalent.

A parameter $\theta\ge-1$ puts the reversion extremum ($\theta=1$) and
superexponentially weighted plane trees ($\theta=-1$) into one theorem.
The proof reduces an exact Lagrange sum to finitely many small-part
sizes and evaluates it by a Poisson change of measure. It gives a
finite-head universality law, all critical windows
$s=1/r+t/\log n$, and explicit all-order centering parameters for the
asymptotically independent small-part counts. Exact algebraic checks and
non-certified numerical diagnostics accompany the proofs.
\end{abstract}
\vfill
\small
\textbf{Status.} Conventional mathematical proofs, not Lean-verified or
peer-reviewed. The identified repository gap is resolved under the
hypotheses stated here; no claim of worldwide publication priority or
resolution of an unrelated longstanding conjecture is made.
\end{titlepage}
\hypersetup{pageanchor=true}
\setcounter{tocdepth}{1}
\tableofcontents
\clearpage

\section{The research target and its boundaries}\label{sec:target}

\subsection{A specific missing asymptotic, not another type theorem}
For a formal power series $f(z)=z+O(z^2)$, let $f\inv$ denote its
compositional inverse. Consider the fixed coefficient ball
\begin{equation}
 \mathcal B_{s,A,C}
 =\left\{z+\sum_{n\ge2}f_nz^n:
       |f_n|\le C A^{n-1}((n-1)!)^s\right\},
 \qquad s,A,C>0.
 \label{eq:basicball}
\end{equation}
Preservation of an $n$th-root coefficient type under inversion does not
specify the amplification at these exact constants $A$ and $C$.
Subexponential factors disappear under $n$th roots, but can still
become arbitrarily large.

The inspected ProveIt package \emph{Sharp Subexponential Reversion}
\cite{subexp} identifies a single extremizer for \eqref{eq:basicball} and
proves its sharp leading logarithm. It gives multiplicative equivalents
above order $1/2$ and at $1/2$, but explicitly leaves the corresponding
answer below $1/2$ open within that article. This is the target of the
present work. The earlier exact weighted-type theorem \cite{weighted}
is not claimed as a new contribution here.

The repository snapshot used for the inspection is
\begin{center}\snapshot.\end{center}
The relevant directories lie under
\begin{quote}\small\ttfamily
Analysis/Transseries/docs/series-and-transseries/.
\end{quote}
The package names in the bibliography, source headings, and this
snapshot are more durable locators than mutable source line numbers.
The source review was targeted: the relevant READMEs, the predecessor's
extremal and localization arguments, and searches for overlapping
reversion results were inspected. This was not an audit of the entire
repository.

\subsection{What is proved here}
The principal theorem supplies, for every fixed $s>0$, a finite list of
explicit constants that determines the extremal inverse coefficient up
to a relative error tending to zero. For example, the previously
uncovered order $s=1/3$ has the formula
\begin{align}
 \log R_n={}&Cn^{2/3}
 +\left(2^{1/3}C+\frac23C^2\right)n^{1/3}\notag\\
 &+6^{1/3}C+\frac53\,2^{1/3}C^2+\frac79C^3+o(1),
 \label{eq:introthird}
\end{align}
where $R_n$ is the supremum of the normalized inverse coefficient over
\eqref{eq:basicball}. The formula has an actual constant term: omitting
it would give an incorrect multiplicative equivalent.

There are three additional consequences. First, only the initial
$\lfloor1/s\rfloor$ weights affect the normalized multiplicative
asymptotic. Second, every reciprocal integer $1/r$ has an explicitly
resolved $1/\log n$ transition window. Third, the same finite implicit
equation computes all centering corrections needed for Gaussian limits
of the accompanying small-part counts. A one-parameter family makes
these statements apply to both extremal reversion and weighted trees.

\subsection{Published precedents and novelty qualifications}
Formal Lagrange inversion is classical; we use its coefficient form as
presented, for example, by Gessel \cite{gessel}. Factorially divergent
coefficient algebras and their composition and inversion rules have
substantial antecedents, including Borinsky \cite{borinsky}.

For superexponential tree weights, Janson, Jonsson and Stef\'ansson
\cite{jjs} already prove condensation, Gaussian and critical Poisson
limits, and total-variation approximation by independent Poisson
variables. Their Remark~2.8 gives the first two exponent terms and
points toward higher corrections. Those probabilistic phenomena are
not discoveries of this article. Our contribution is the explicit
all-order pressure and centering formula, its reversion application,
and the finite-head and transition-window consequences proved below.

The principal proofs are self-contained after formal Lagrange inversion
and elementary facts about factorials and Poisson variables. In
particular, no unreviewed asymptotic theorem from the repository is
needed as a black-box input. The literature check is not an exhaustive
priority search; the manuscript should receive independent mathematical
review before its conclusions are treated as established results.

\begin{table}[htbp]
\centering\small
\begin{tabular}{p{.32\textwidth}p{.61\textwidth}}
\toprule
Statement & Scope and status\\
\midrule
Exact extremizer & Known in the predecessor; short proof included.\\
All-order exponent coefficients & Explicit formula and proof for every fixed $s>0$.\\
Complete multiplicative equivalent & Proved for eventually shifted-factorial tails and arbitrary positive finite heads.\\
All reciprocal critical windows & Proved, uniformly on bounded window parameters.\\
Poisson small-part law & Generalized and given explicit centers; the tree phenomenon has published precedents.\\
Numerical evidence & Finite exact algebra checks plus floating-point diagnostics; not an asymptotic proof.\\
Beyond the theorem's hypotheses & General oscillatory tails, cancellation, and the joint limit $s\downarrow0$ remain separate questions.\\
\bottomrule
\end{tabular}
\caption{Claim and dependency ledger.}
\end{table}

\section{The main theorem}\label{sec:main}

\subsection{Weights and the interpolation family}
\begin{definition}[Eventually shifted-factorial weights]\label[definition]{def:weights}
Fix $s>0$. A positive sequence $(w_j)_{j\ge1}$ is admissible here if
there exist $D_*>0$, $\nu\ge0$, and an integer $j_0\ge1$ such that
\begin{equation}
 w_j=D_*\,\Gamma(j+\nu+1)^s\qquad(j\ge j_0).
 \label{eq:tail}
\end{equation}
The finitely many weights $w_1,\ldots,w_{j_0-1}$ may be arbitrary
positive numbers. In particular, log-convexity of the entire sequence is
not required.
\end{definition}

The weights in the predecessor are the special cases
\begin{equation}
 w_j=\left(\frac{\Gamma(j+\nu+1)}{\Gamma(\nu+2)}\right)^s,
 \qquad \nu\in\{0,1\}.
 \label{eq:repoweights}
\end{equation}
They satisfy $w_1=1$. The shift $\nu=0$ gives $w_j=(j!)^s$;
$\nu=1$ gives $w_j=((j+1)!/2)^s$.

Fix $C>0$ and $\theta\ge-1$. All the following definitions involving
an infinite series are formal. Put
\begin{equation}
 W(z)=\sum_{j\ge1}w_jz^j,
 \qquad
 F_\theta(z)=
 \begin{cases}
 z(1-\theta CW(z))^{1/\theta},&\theta\ne0,\\
 z\exp(-CW(z)),&\theta=0,
 \end{cases}
 \label{eq:Ftheta}
\end{equation}
and define
\begin{equation}
 G_\theta=F_\theta\inv,
 \qquad R_n^{(\theta)}=
 \frac{[z^n]G_\theta(z)}{Cw_{n-1}},\qquad n\ge2.
 \label{eq:Rtheta}
\end{equation}
The linear coefficient of $F_\theta$ is $1$, so the inverse exists
uniquely in $\R[[z]]$. We shall prove that $R_n^{(\theta)}>0$.
The restriction $\theta\ge-1$ is a positivity condition on finite
Lagrange coefficients, not an analyticity assumption on $W$.

At $\theta=1$, $F_1=z-C\sum_{j\ge1}w_jz^{j+1}$ is the extremal
negative series. At $\theta=-1$, $G_{-1}$ satisfies
\begin{equation}
 G_{-1}(z)=z\bigl(1+CW(G_{-1}(z))\bigr),
 \label{eq:treeeq}
\end{equation}
the formal equation for plane rooted trees with leaf weight $1$ and
weight $Cw_j$ for a vertex of outdegree $j\ge1$.

\subsection{A finite polynomial determines the answer}
Choose an integer $J\ge1$ such that
\begin{equation}
 (J+1)s>1.
 \label{eq:J}
\end{equation}
Define the finite polynomial and its rational transform
\begin{equation}
 a_j=Cw_j,\qquad
 A(t)=\sum_{j=1}^J a_jt^j,
 \qquad B(t)=\frac{tA'(t)}{1-\theta A(t)}.
 \label{eq:AB}
\end{equation}
Unlike $W$, these expressions are analytic near the origin.
For any real $a$, let
\begin{equation}
 \Psi_a(u)=\int_0^u(1-v)^{-a}\,dv
 =\sum_{r\ge1}\frac{\rise{a}{r-1}}{r!}u^r,
 \label{eq:Psi}
\end{equation}
where the rising factorial of length zero is $1$. Thus
$\Psi_1(u)=-\log(1-u)$; for $a\ne1$,
$\Psi_a(u)=((1-u)^{1-a}-1)/(a-1)$.

\begin{theorem}[Finite-jet pressure law]\label{thm:main}
Under \cref{def:weights}, fix $C>0$, $\theta\ge-1$, and $J$ satisfying
\eqref{eq:J}. Define
\begin{equation}
 p_k=\frac1k[t^k]\,\Psi_{ks}(B(t)),\qquad 1\le k\le J.
 \label{eq:pk}
\end{equation}
Then
\begin{equation}
 \boxed{\displaystyle
 \log R_n^{(\theta)}=\sum_{k=1}^Jp_kn^{1-ks}+o(1).}
 \label{eq:mainlog}
\end{equation}
Equivalently,
\begin{equation}
 R_n^{(\theta)}=
 \exp\left\{\sum_{k=1}^Jp_kn^{1-ks}\right\}(1+o(1)).
 \label{eq:mainmult}
\end{equation}
For a fixed $s$, terms with $ks>1$ may be omitted. The coefficient $p_k$
depends only on $s,\theta,a_1,\ldots,a_k$, not on the rest of the
weight sequence or the chosen truncation order $J\ge k$.
\end{theorem}

In particular, when $0<s\le1$, taking $J=\lfloor1/s\rfloor$ supplies
exactly the terms that do not tend to zero. For $s>1$, taking $J=1$
shows $R_n^{(\theta)}\to1$. At $s=1$ the limit is $\exp(Cw_1)$,
independent of $\theta$.

Here and throughout, a complete multiplicative equivalent means a
formula with relative error $o(1)$. It does not mean that every vanishing
correction after that relative error has been calculated.

\begin{theorem}[Finite analytic pressure]\label{thm:pressure}
With the same hypotheses, let $y=y(x)$ be the unique analytic germ
with $y(0)=0$ satisfying
\begin{equation}
 y=x(1-B(y))^{-s}.
 \label{eq:y}
\end{equation}
Set $d(x)=B(y(x))$ and
\begin{equation}
 P(x)=L_\theta(A(y(x)))
       +s\bigl(\log(1-d(x))+d(x)\bigr),
 \label{eq:pressure}
\end{equation}
where
\begin{equation}
 L_\theta(v)=
 \begin{cases}
 -\theta^{-1}\log(1-\theta v),&\theta\ne0,\\
 v,&\theta=0.
 \end{cases}
 \label{eq:Ltheta}
\end{equation}
For sufficiently large $n$, this real germ is defined at $x=n^{-s}$,
and
\begin{equation}
 R_n^{(\theta)}=\exp\{nP(n^{-s})\}(1+o(1)),
 \qquad xP'(x)=d(x).
 \label{eq:resummed}
\end{equation}
Its Taylor coefficients through order $J$ are precisely \eqref{eq:pk}.
\end{theorem}

The pressure in \eqref{eq:pressure} uses a \emph{finite} polynomial.
There is no substitution of a positive real number into the divergent
series $W$. This distinction is necessary: replacing $A$ by $CW$ in an
analytic saddle equation would not be justified.

\subsection{Uniformity needed later}
The conclusions of \cref{thm:main,thm:pressure} hold uniformly when
$s,C$, and $\theta$ vary in compact subsets of
$(0,\infty)$, $(0,\infty)$, and $[-1,\infty)$, respectively, provided
one fixed $J$ has $(J+1)s_{\min}>1$ and the tail description is uniform.
For example, one may fix $j_0$ and $\nu$, let $D_*$ and the finitely many
exceptional weights be positive continuous functions of the parameters,
and bound them above and away from zero on those compact sets.
This version includes \eqref{eq:repoweights} with varying $s$.
The proof below identifies the uniform estimates explicitly.

\section{Exact Lagrange coefficients and the extremum}\label{sec:exact}

\subsection{The positive finite sum}
Let $m=n-1$, and write
\begin{equation}
 P_w(m,k)=\sum_{\substack{j_1+\cdots+j_k=m\\j_i\ge1}}
                   \prod_{i=1}^k w_{j_i}.
 \label{eq:Pw}
\end{equation}
This is a finite sum over ordered positive compositions of $m$.

\begin{proposition}[Exact coefficient formula]\label[proposition]{prop:lagrange}
For $n\ge2$,
\begin{equation}
 R_n^{(\theta)}=
 \frac1{nCw_m}\sum_{k=1}^m
 \frac{\prod_{r=0}^{k-1}(n+\theta r)}{k!}\,C^kP_w(m,k).
 \label{eq:exactR}
\end{equation}
Every summand is positive and the $k=1$ summand equals $1$.
\end{proposition}
\begin{proof}
Formal Lagrange inversion \cite{gessel} gives
\[
 [z^n]G_\theta=
 \frac1n[t^{n-1}](1-\theta CW(t))^{-n/\theta}
\]
for $\theta\ne0$, with $\exp(nCW(t))$ at $\theta=0$.
The coefficient of $(CW)^k$ is
$\prod_{r=0}^{k-1}(n+\theta r)/k!$, including the continuous
$\theta=0$ interpretation. Since $k\le n-1$ and $\theta\ge-1$,
all factors are strictly positive. Expansion of $W^k$ gives
\eqref{eq:exactR}; its first term is $1$ after normalization.
\end{proof}

No convergence is used in this argument. A coefficient of degree $m$
in a power of $W$ depends on only $w_1,\ldots,w_m$.

\begin{proposition}[Simultaneous extremality]\label[proposition]{prop:extreme}
For arbitrary positive weights and $A_0,C>0$, let
\[
 \mathcal B(A_0,C;w)=
 \left\{z+\sum_{j\ge1}f_{j+1}z^{j+1}:\
       |f_{j+1}|\le C A_0^jw_j\right\}.
\]
The single series
\begin{equation}
 F_*(z)=z-C\sum_{j\ge1}A_0^jw_jz^{j+1}
 \label{eq:extremizer}
\end{equation}
attains
\begin{equation}
 \sup_{f\in\mathcal B(A_0,C;w)}
 \frac{|[z^n]f\inv|}{CA_0^{n-1}w_{n-1}}
 =R_n^{(1)}
 \label{eq:extremalR}
\end{equation}
for every $n\ge2$ simultaneously.
\end{proposition}
\begin{proof}
For $f=z(1+H)$, Lagrange inversion expresses the inverse coefficient as
$\frac1n[z^{n-1}](1+H)^{-n}$.
The coefficients of $(1-u)^{-n}$ are nonnegative. Taking absolute
values and majorizing each coefficient of $H$ by $CA_0^jw_j$
therefore gives the coefficient of $F_*\inv$.
For $F_*$ every contribution has the same positive sign, so equality
holds. Finally, replacing $z$ by $A_0z$ contributes the common factor
$A_0^{n-1}$. This proves both the extremum and its independence of
$A_0$ after normalization.
\end{proof}

This proposition is an exact statement about a coefficient ball, not
about the inverse of an arbitrary positive or signed series. The
asymptotic conclusions about this supremum do not assert that every
member of the ball attains it asymptotically.

\section{Localization at one large part}\label{sec:localization}

The first reduction is from all compositions to those with a unique
part near $m=n-1$. The argument also records why a fixed finite head
of the weights does not obstruct the analysis.

\begin{lemma}[Convolution bound]\label[lemma]{lem:convolution}
Under \cref{def:weights}, there is a constant $K\ge1$ such that
\begin{equation}
 \sum_{j=1}^{m-1}w_jw_{m-j}\le K w_{m-1}\quad(m\ge2),
 \qquad
 P_w(m,k)\le K^{k-1}w_{m-k+1}.
 \label{eq:conv}
\end{equation}
\end{lemma}
\begin{proof}
There are constants $c_1,c_2>0$ such that, for every $j\ge1$,
\begin{equation}
 c_1(j!)^s(j+1)^{s\nu}\le w_j
 \le c_2(j!)^s(j+1)^{s\nu}.
 \label{eq:gammacomparison}
\end{equation}
This follows from the gamma quotient asymptotic and adjustment of
finitely many constants; for integer $\nu$ it follows directly from a
finite product.

For large $m$, and $1\le j\le m/2$, both $m-j$ and $m-1$ are in the
exact tail. Their quotient is
\[
 \frac{w_{m-j}}{w_{m-1}}
 =\left(\prod_{r=1}^{j-1}(m-j+\nu+r)\right)^{-s}
 \le j^{-s(j-1)}.
\]
Consequently
\[
 \frac{w_jw_{m-j}}{w_{m-1}}
 \le c_2(j!)^s(j+1)^{s\nu}j^{-s(j-1)}.
\]
Stirling's upper bound makes the right side at most a constant times
$j^{s\nu+3s/2}\ee^{-sj}$. This is summable. Reflection about $m/2$
and an increase of $K$ for the remaining finitely many $m$ prove the
first assertion.

For the second, induct on $k$. Splitting off the first part and using
the induction hypothesis gives
\[
 P_w(m,k)\le K^{k-2}
 \sum_{j=1}^{m-k+1}w_jw_{m-k+2-j}
 \le K^{k-1}w_{m-k+1}.
\]
The case $k=1$ is an equality.
\end{proof}

Let $T_{n,k}$ be the $k$th summand in \eqref{eq:exactR}.
The elementary estimate
$\prod_{r=0}^{k-1}(n+\theta r)\le(c_\theta n)^k$,
where $c_\theta=1+\max(\theta,0)$, combines with
\cref{lem:convolution} to give constants $c,D>0$ such that
\begin{equation}
 T_{n,k}\le c\frac{(D n^{1-s})^{k-1}}{k!}
 \qquad(1\le k\le n/2)
 \label{eq:lengthtail}
\end{equation}
for large $n$. Indeed, the quotient $w_{n-k}/w_{n-1}$ consists of
$k-1$ inverse $s$th powers of factors at least $n/2$.
For $k>n/2$, monotonicity in the exact tail, with a fixed adjustment for
the finite head, gives
\begin{equation}
 \sum_{k>n/2}T_{n,k}
 \le \exp\{-\tfrac{s}{2}n\log n+O(n)\}.
 \label{eq:longlength}
\end{equation}
To check the entropy factor in this last bound, note that
$(c_\theta n)^k/k!$ is $\exp(O(n))$ when $k>n/2$;
$C^{k-1}K^{k-1}$ has the same bound. The remaining weight quotient has
logarithm at most $-\frac{s}{2}n\log n+O(n)$.

\begin{lemma}[One-large-part localization]\label[lemma]{lem:localization}
Choose real exponents satisfying
\begin{equation}
 \max(1-s,0)<\kappa<\alpha<1,
 \qquad h=\lfloor n^\alpha\rfloor.
 \label{eq:cutoff}
\end{equation}
For large $n$, the total contribution to \eqref{eq:exactR} from
compositions with no part at least $m-h$ is $o(1)$.
More precisely, it is bounded by a sum of terms of the forms
$\exp(-c n^\alpha\log n)$ and
$\exp(-c' n^\kappa\log n)$, with positive constants.
\end{lemma}
\begin{proof}
First restrict the number of parts to $k\le t=\lfloor n^\kappa\rfloor$.
Since $h<m/2$ eventually, concentration of a factorial product under
an upper bound $m-h$ gives
\begin{equation}
 \prod_{i=1}^k j_i!\le(m-h)!h!.
 \label{eq:cap}
\end{equation}
One proof permits zero parts and moves a unit from a smaller positive
part to a larger part that is below the cap. Such a move does not
decrease the factorial product. At a maximum at most two parts are
nonzero, and they have sizes $m-h$ and $h$.

By \eqref{eq:gammacomparison}, the normalized weight product is at most
\[
 c_1^{-1}c_2^k\binom{m}{h}^{-s}
       \prod_{i=1}^k(j_i+1)^{s\nu}
 \le c_1^{-1}c_2^k\binom{m}{h}^{-s}(1+m/k)^{s\nu k}.
\]
We have discarded a denominator $(m+1)^{s\nu}\ge1$.
There are $\binom{m-1}{k-1}$ positive compositions of length $k$.
The logarithms of this count, of the Lagrange factor, and of all fixed
constants together are $O(k\log(\ee n/k)+\log n)$.
The total for $k\le t$ is thus bounded by
\begin{equation}
 \exp\left\{-s\log\binom{m}{h}
       +O\left(t\log\frac{\ee n}{t}+\log n\right)\right\}.
 \label{eq:capentropy}
\end{equation}
Here
$\log\binom{m}{h}=(1-\alpha+o(1))n^\alpha\log n$,
which dominates the positive error because $\alpha>\kappa$.

For $k>t$, use \eqref{eq:lengthtail} and \eqref{eq:longlength}.
The ratio $t/n^{1-s}$ tends to infinity. The inequality
$k!\ge(k/\ee)^k$, followed by a geometric tail estimate, gives
\[
 \sum_{t<k\le n/2}T_{n,k}
 \le\exp\left\{-t\log\frac{t}{\ee D n^{1-s}}+O(\log n)\right\}.
\]
Its exponent is a negative constant multiple of
$n^\kappa\log n$. These bounds prove the claim.
\end{proof}

The error in this lemma is \emph{absolute} in the normalized quantity
$R_n^{(\theta)}$. Since $R_n^{(\theta)}\ge1$, an absolute $o(1)$ is
also harmless for a multiplicative equivalent. In the next reduction,
where sums can be exponentially large on a sublinear scale, a relative
estimate will be needed instead.

\section{The exact cloud and the finite-size cutoff}\label{sec:cloud}

\subsection{Grouping the small parts}
In a localized composition the part of size at least $m-h$ is unique.
Call it the large part and write its size as $m-d$. If $m_j$ counts the
remaining parts of size $j$, then
\begin{equation}
 \ell=\sum_{j\ge1}m_j,\qquad d=\sum_{j\ge1}jm_j\le h.
 \label{eq:ellandd}
\end{equation}
There are $(\ell+1)!/\prod_jm_j!$ ordered compositions with this
multiplicity vector. Multiplication by the $k=\ell+1$ Lagrange factor
in \eqref{eq:exactR} cancels both that factorial and the first factor
$n$. We obtain the exact formula
\begin{equation}
 R_{n,\mathrm{loc}}^{(\theta)}(h)=
 \sum_{\substack{m_1,m_2,\ldots\ge0\\\sum jm_j\le h}}
 \frac{\prod_{r=1}^{\ell}(n+\theta r)}{\prod_jm_j!}
 \frac{w_{m-d}}{w_m}\prod_{j\ge1}a_j^{m_j}.
 \label{eq:cloudexact}
\end{equation}
The all-zero multiplicity vector contributes $1$. In this section
$a_j=Cw_j$ is defined for all $j\ge1$, not only $j\le J$.

\begin{lemma}[Only finitely many small sizes matter]\label[lemma]{lem:finitecut}
Let $S_{n,J}(h)$ denote the sum in \eqref{eq:cloudexact} restricted to
$m_j=0$ for all $j>J$. Under \eqref{eq:J} and \eqref{eq:cutoff},
\begin{equation}
 R_n^{(\theta)}=S_{n,J}(h)(1+o(1)).
 \label{eq:finitecut}
\end{equation}
\end{lemma}
\begin{proof}
All summands are nonnegative. Compare a summand with one obtained by
adding a part of size $j$ while retaining $d+j\le h$. The ratio is at
most
\begin{equation}
 \frac{\lambda_j^+}{m_j+1},\qquad
 \lambda_j^+=(n+\theta_+h)a_j(n+\nu-h)^{-sj},
 \qquad\theta_+=\max(\theta,0).
 \label{eq:insertion}
\end{equation}
For the first factor use $n+\theta(\ell+1)\le n+\theta_+h$.
For the large-part quotient, all relevant indices are eventually in
the exact tail and each of the $j$ denominator factors is at least
$n+\nu-h$. This proves \eqref{eq:insertion}.

For $j$ in the exact tail and $j<h$, successive intensities have ratio
\[
 \frac{\lambda_{j+1}^+}{\lambda_j^+}
 =\left(\frac{j+\nu+1}{n+\nu-h}\right)^s=o(1)
\]
uniformly in that range. At the finitely many exceptional indices the
same conclusion follows from a fixed weight quotient times $O(n^{-s})$.
Hence
\begin{equation}
 \sum_{j=J+1}^h\lambda_j^+
 =O\bigl(n^{1-(J+1)s}\bigr)=o(1).
 \label{eq:insertionsmall}
\end{equation}

Delete every part of size greater than $J$ from a cloud. Its remaining
low-size cloud is still admissible. Reinsert the deleted parts in a
fixed order and apply \eqref{eq:insertion} repeatedly. Summing their
multiplicities bounds the ratio to the original low-size summand by
$\prod_{j>J}\sum_{r\ge0}(\lambda_j^+)^r/r!$.
Extending the sum beyond the original size cutoff only increases it.
Thus
\[
 S_{n,J}(h)\le R_{n,\mathrm{loc}}^{(\theta)}(h)
 \le S_{n,J}(h)\exp\left(\sum_{j>J}\lambda_j^+\right)
 =S_{n,J}(h)(1+o(1)).
\]
Finally, \cref{lem:localization} adds only $o(1)$, and
$S_{n,J}(h)\ge1$. This proves \eqref{eq:finitecut}.
\end{proof}

This is the source of finite-jet universality. A size $j$ has a
natural intensity of order $n^{1-js}$. Once this tends to zero, the
total effect of all such omitted sizes is multiplicatively negligible.
The proof justifies this intuition without evaluating the divergent
weight generating function.

\section{The interaction and its Poisson evaluation}\label{sec:poissonproof}

\subsection{An exact product becomes a smooth local interaction}
For $u$ near zero define
\begin{equation}
 E_\theta(u)=
 \begin{cases}
 \bigl((1+\theta u)\log(1+\theta u)-\theta u\bigr)/\theta,
    &\theta\ne0,\\
 0,&\theta=0,
 \end{cases}
 \label{eq:Etheta}
\end{equation}
and, for $v$ near zero, define
\begin{equation}
 b(v)=(1-v)\log(1-v)+v,
 \qquad \mathcal H(u,v)=E_\theta(u)+sb(v).
 \label{eq:interaction}
\end{equation}
Its derivatives are
\begin{align}
 \mathcal H_u(u,v)&=\log(1+\theta u),&
 \mathcal H_v(u,v)&=-s\log(1-v),\label{eq:gradH}\\
 \mathcal H_{uu}(u,v)&=\frac{\theta}{1+\theta u},&
 \mathcal H_{vv}(u,v)&=\frac{s}{1-v},&
 \mathcal H_{uv}&=0.\label{eq:hessianH}
\end{align}
The Hessian need not be positive: the tree endpoint has $\theta=-1$.
Only boundedness of the Hessian near the origin is used.

\begin{lemma}[Uniform product approximation]\label[lemma]{lem:product}
Uniformly for integer $0\le\ell\le d\le h=o(n)$,
\begin{equation}
 \prod_{r=1}^\ell(n+\theta r)\frac{w_{n-1-d}}{w_{n-1}}
 =n^{\ell-sd}
 \exp\left\{n\mathcal H(\ell/n,d/n)+O(h/n)\right\}.
 \label{eq:productapprox}
\end{equation}
\end{lemma}
\begin{proof}
The exact tail identity gives
\[
 \frac{w_{n-1-d}}{w_{n-1}}
 =\prod_{r=1}^d(n+\nu-r)^{-s}.
\]
After removing $n^{\ell-sd}$, the logarithm of the left side in
\eqref{eq:productapprox} is
\[
 \sum_{r=1}^\ell\log(1+\theta r/n)
 -s\sum_{r=1}^d\log(1+(\nu-r)/n).
\]
The first sum is its corresponding integral plus $O(\ell/n)$.
Replacing $1+(\nu-r)/n$ by $1-r/n$ changes the second by $O(d/n)$,
and its sum-to-integral error is also $O(d/n)$.
The integrals are $nE_\theta(\ell/n)$ and $nsb(d/n)$.
All derivatives involved are uniformly bounded because $h/n\to0$.
\end{proof}

Set $x=n^{-s}$. By \cref{lem:product},
\begin{equation}
 S_{n,J}(h)=(1+o(1))\!
 \sum_{\substack{m_1,\ldots,m_J\ge0\\d=\sum jm_j\le h}}
 \prod_{j=1}^J\frac{(na_jx^j)^{m_j}}{m_j!}
       \exp\{n\mathcal H(\ell/n,d/n)\}.
 \label{eq:interactionsum}
\end{equation}
The relative error is uniform over every summand. This is a
finite-dimensional sum, even though the original weight series is
divergent.

\subsection{A small saddle, defined without global maximization}
The implicit function theorem applied to
$y-x(1-B(y))^{-s}$ gives the unique analytic germ in \eqref{eq:y}.
For sufficiently small positive $x$, $y>0$, $1-\theta A(y)>0$, and
$1-B(y)>0$. Define
\begin{equation}
 \beta_j(x)=\frac{a_jy(x)^j}{1-\theta A(y(x))},
 \quad u_*(x)=\sum_{j=1}^J\beta_j(x),
 \quad v_*(x)=\sum_{j=1}^Jj\beta_j(x).
 \label{eq:betas}
\end{equation}
Then
\begin{equation}
 u_* =\frac{A(y)}{1-\theta A(y)},\qquad
 v_*=B(y),\qquad
 1+\theta u_*=(1-\theta A(y))^{-1}.
 \label{eq:uv}
\end{equation}
Equation \eqref{eq:y} is equivalent to the stationary identities
\begin{equation}
 \log\frac{\beta_j}{a_jx^j}
 =\mathcal H_u(u_*,v_*)+j\mathcal H_v(u_*,v_*),
 \qquad1\le j\le J.
 \label{eq:stationarity}
\end{equation}
We have $\beta_j=a_jx^j(1+O(x))$ and $u_*,v_*=O(x)$.
No assertion that this stationary point is a global maximum is needed.
The next lemma evaluates its neighborhood and bounds the entire
remaining cutoff region at the same time.

\subsection{A cutoff exponential-integrability lemma}
\begin{lemma}[Poisson tilt remainder]\label[lemma]{lem:Poissontilt}
Fix $J$. Let $N_1,\ldots,N_J$ be independent Poisson variables with
means $\mu_j\ge0$, let
\[
 L=\sum_jN_j,\quad D=\sum_jjN_j,\quad
 v=\sum_j\mu_j,
\]
and suppose $v=o(n)$, $h=o(n)$, $h\to\infty$, and $v/h\to0$.
Write
\[
 U=\max(|L-\E L|,|D-\E D|).
\]
For every fixed $K_0>0$,
\begin{equation}
 \E\left[(\exp(K_0U^2/n)-1)\ind_{D\le h}\right]=o(1),
 \qquad \PP(D>h)=o(1).
 \label{eq:Poissonexp}
\end{equation}
Consequently, if a random remainder $\mathcal E_n$ on $\{D\le h\}$
satisfies $|\mathcal E_n|\le K_0U^2/n$, then
\begin{equation}
 \E[\exp(\mathcal E_n)\ind_{D\le h}]=1+o(1),
 \qquad
 \E\left|\exp(\mathcal E_n)\ind_{D\le h}-1\right|=o(1).
 \label{eq:L1tilt}
\end{equation}
\end{lemma}
\begin{proof}
For a centered weighted Poisson sum
$X=\sum c_j(N_j-\mu_j)$ with $0\le c_j\le J$, its cumulant generating
function is
\[
 \log\E\ee^{tX}=\sum_j\mu_j(\ee^{tc_j}-1-tc_j).
\]
For $|t|\le1/(2J)$, Taylor's inequality bounds this by a constant
multiple of $v t^2$. Applying Chernoff's bound with
$t$ proportional to $u/(v+u)$, truncated to that interval, gives
constants depending only on $J$ such that
\begin{equation}
 \PP(U>u)\le4\exp\left\{-c\frac{u^2}{v+u}\right\}
 \quad(u>0).
 \label{eq:Bernstein}
\end{equation}
When $v=0$ the variables vanish and the conclusion is immediate; one
may otherwise replace $v$ by $v+1$ in this bound.

On $D\le h$, both $L\le D$ and $\E D\le Jv=o(h)$ imply $U\le2h$
for all sufficiently large $n$. Tail integration therefore bounds the
first expectation in \eqref{eq:Poissonexp} by
\begin{equation}
 \int_0^{2h}\frac{2K_0u}{n}\ee^{K_0u^2/n}\PP(U>u)\,du.
 \label{eq:tailintegral}
\end{equation}
Because $(v+1+h)/n\to0$, the positive exponent in this integral is
at most half the magnitude of the negative exponent in
\eqref{eq:Bernstein}, uniformly for $0\le u\le2h$.
On $u\le v+1$ the resulting integral is $O((v+1)/n)$ by a Gaussian
integral. On $u>v+1$ it is $O(1/n)$ after using an exponential tail
$\exp(-c'u)$. Both tend to zero.

Markov's inequality gives $\PP(D>h)\le Jv/h=o(1)$.
Finally, $|\ee^z-1|\le\ee^{|z|}-1$ for real $z$.
Combining these facts yields \eqref{eq:L1tilt}.
\end{proof}

\begin{remark}[Why the cutoff is retained]
One must not replace the truncated expectation by
$\E\exp(cN^2/n)$ for a nontrivial Poisson variable $N$ and $c>0$;
that expectation is infinite. The condition $D\le h=o(n)$ is part of
the proof, not a dispensable convenience. \Cref{lem:Poissontilt} controls
the whole required region without making that invalid extension.
\end{remark}

\subsection{Evaluation by exponential change of measure}
Take independent Poisson variables with means
\begin{equation}
 \mu_{j,n}=n\beta_j(n^{-s}),\qquad1\le j\le J.
 \label{eq:means}
\end{equation}
Their total mean is $O(n^{1-s})$, or tends to zero when $s>1$.
In either case it is $o(h)$ for the cutoff chosen in
\eqref{eq:cutoff}.

Use \eqref{eq:stationarity} to change the factorial weights in
\eqref{eq:interactionsum} from means $na_jx^j$ to means $n\beta_j$.
If $L=\sum N_j$ and $D=\sum jN_j$, the result is
\begin{equation}
 S_{n,J}(h)=(1+o(1))\exp\{n\widetilde P(x)\}
     \E[\exp(\mathcal E_n)\ind_{D\le h}],
 \label{eq:tiltexact}
\end{equation}
where
\begin{align}
 \widetilde P(x)
 &=u_*+\mathcal H(u_*,v_*)
   -u_*\mathcal H_u(u_*,v_*)-v_*\mathcal H_v(u_*,v_*),
 \label{eq:Pvariational}\\
 \mathcal E_n
 &=n\bigl[\mathcal H(L/n,D/n)-\mathcal H(u_*,v_*)\notag\\
 &\hspace{17mm}
 -\mathcal H_u(u_*,v_*)(L/n-u_*)
 -\mathcal H_v(u_*,v_*)(D/n-v_*)\bigr].
 \label{eq:Eremainder}
\end{align}
The Hessian in \eqref{eq:hessianH} is bounded on all line segments
involved when $D\le h$. Taylor's theorem gives
\begin{equation}
 |\mathcal E_n|\le\frac{K_0}{n}
 \bigl((L-nu_*)^2+(D-nv_*)^2\bigr).
 \label{eq:remainderbound}
\end{equation}
Thus \cref{lem:Poissontilt} makes the expectation in
\eqref{eq:tiltexact} equal to $1+o(1)$.

Substituting \eqref{eq:uv} and \eqref{eq:gradH} into
\eqref{eq:Pvariational} simplifies it to
\[
 \widetilde P(x)=
 -\theta^{-1}\log(1-\theta A(y))
 +s\bigl(\log(1-B(y))+B(y)\bigr)
\]
for $\theta\ne0$, with the stated continuous limit at zero. This is
exactly $P(x)$ in \eqref{eq:pressure}. Together with
\cref{lem:finitecut}, we have proved the multiplicative pressure
formula in \cref{thm:pressure}.

\section{Extracting the pressure coefficients}\label{sec:extraction}

\subsection{The envelope identity}
For positive $x$ and positive $\beta=(\beta_1,\ldots,\beta_J)$ near
the small stationary point, introduce
\begin{equation}
 \mathcal V(\beta,x)=
 \sum_{j=1}^J\beta_j
       \left(1+\log\frac{a_jx^j}{\beta_j}\right)
 +\mathcal H\left(\sum_j\beta_j,\sum_jj\beta_j\right).
 \label{eq:variational}
\end{equation}
Its partial derivatives with respect to $\beta_j$ vanish at
\eqref{eq:betas}, precisely by \eqref{eq:stationarity}.
At that point its value is $P(x)$, by \eqref{eq:Pvariational}.
Differentiating the composite function therefore gives
\begin{equation}
 xP'(x)=\sum_{j=1}^Jj\beta_j(x)=B(y(x)).
 \label{eq:envelope}
\end{equation}
This differentiation is performed first for $x>0$; both sides are
analytic at zero and hence the identity also holds as an identity of
power series. It uses stationarity, not a global variational principle.

\subsection{Lagrange inversion a second time}
Apply formal Lagrange inversion to $y=x(1-B(y))^{-s}$ and the
function $B$. For $k\ge1$,
\begin{equation}
 [x^k]B(y(x))
 =\frac1k[t^{k-1}]B'(t)(1-B(t))^{-ks}.
 \label{eq:LagrangeB}
\end{equation}
Since $\Psi_{ks}'(u)=(1-u)^{-ks}$, equations
\eqref{eq:envelope} and \eqref{eq:LagrangeB} imply
\begin{align}
 [x^k]P(x)
 &=\frac1{k^2}[t^{k-1}]\frac{d}{dt}\Psi_{ks}(B(t))\notag\\
 &=\frac1k[t^k]\Psi_{ks}(B(t)).
 \label{eq:pkproof}
\end{align}
This proves \eqref{eq:pk}. A coefficient of degree $k$ uses only
$a_1,\ldots,a_k$, proving the finite-jet assertion.

Finally, analyticity gives
\[
 nP(n^{-s})=\sum_{k=1}^Jp_kn^{1-ks}
                 +O\bigl(n^{1-(J+1)s}\bigr).
\]
The remainder tends to zero by \eqref{eq:J}. This completes the proofs
of \cref{thm:main,thm:pressure}.

\subsection{Why the uniform version follows}
On a compact positive $s$-interval, choose once and for all
\[
 \max(1-s_{\min},0)<\kappa<\alpha<1.
\]
The summable majorant in \cref{lem:convolution}, the tail constants,
and the insertion estimates are uniform under the stated bounds on the
finite head and tail. The strict inequality $(J+1)s_{\min}>1$ supplies
a common positive decay exponent for \eqref{eq:insertionsmall}.
The implicit function theorem supplies a common neighborhood of zero
because $J$ is fixed and the finite coefficients range over a compact
set. The Hessian bounds and the Poisson tail constants are also
uniform. Finally, Taylor's theorem for $P$ has a common remainder
constant on a smaller common neighborhood. These observations prove
the uniformity statement following \cref{thm:pressure}.

\section{Explicit coefficients and every critical point}\label{sec:coefficients}

\subsection{The first four universal polynomials}
Expanding \eqref{eq:pk} gives
\begin{align}
 p_1={}&a_1,\label{eq:p1}\\
 p_2={}&a_2+\frac{s+\theta}{2}a_1^2,\label{eq:p2}\\
 p_3={}&a_3+(2s+\theta)a_1a_2
 +\left(\frac{s^2}{2}+s\theta+\frac{s}{6}
                  +\frac{\theta^2}{3}\right)a_1^3,\label{eq:p3}\\
 p_4={}&a_4+(3s+\theta)a_1a_3
       +\left(2s+\frac\theta2\right)a_2^2\notag\\
 &+(4s^2+5s\theta+s+\theta^2)a_1^2a_2\notag\\
 &+\left(\frac{2s^3}{3}+2s^2\theta+\frac{s^2}{2}
       +\frac{3s\theta^2}{2}+\frac{s\theta}{2}
       +\frac{s}{12}+\frac{\theta^3}{4}\right)a_1^4.
 \label{eq:p4}
\end{align}
These polynomials are valid for the whole interpolation family. In
particular, the terms independent of $\theta$ in the higher
coefficients must not be dropped when passing between its endpoints.
The accompanying exact checker verifies these displayed formulas
against the coefficient-extraction formula.

For the reversion case $\theta=1$, they simplify to
\begin{align}
 p_2={}&a_2+\frac{1+s}{2}a_1^2,\label{eq:revp2}\\
 p_3={}&a_3+(1+2s)a_1a_2
       +\frac{3s^2+7s+2}{6}a_1^3,\label{eq:revp3}\\
 p_4={}&a_4+(1+3s)a_1a_3
       +\left(\frac12+2s\right)a_2^2
       +(4s^2+6s+1)a_1^2a_2\notag\\
 &+\left(\frac{2s^3}{3}+\frac{5s^2}{2}
                         +\frac{25s}{12}+\frac14\right)a_1^4.
 \label{eq:revp4}
\end{align}

\subsection{Reciprocal-integer constants}
\begin{corollary}[Closed critical coefficient]\label[corollary]{cor:critical}
At $s=1/r$, the coefficient of the constant term $n^{1-rs}$ is
\begin{equation}
 \boxed{\displaystyle
 p_r=\frac1r[t^r]\log
 \frac{1-\theta A(t)}{1-\theta A(t)-tA'(t)}.}
 \label{eq:criticalcoefficient}
\end{equation}
Its coefficient of the monomial $a_1^r$ is
\begin{equation}
 [a_1^r]p_r=\frac{(\theta+1)^r-\theta^r}{r^2}.
 \label{eq:purecritical}
\end{equation}
\end{corollary}
\begin{proof}
At $rs=1$, $\Psi_{rs}(u)=-\log(1-u)$. Substitution of
$B=tA'/(1-\theta A)$ into \eqref{eq:pk} proves
\eqref{eq:criticalcoefficient}. To extract the pure $a_1^r$ coefficient,
set $A(t)=a_1t$ and expand
\[
 \frac1r\left(\log(1-\theta a_1t)
       -\log(1-(\theta+1)a_1t)\right).
\]
The coefficient of $t^r$ is \eqref{eq:purecritical}.
\end{proof}

At the reversion endpoint this coefficient is $(2^r-1)/r^2$.
At the tree endpoint it is $(-1)^{r+1}/r^2$.
The latter alternating sign does not contradict positivity of
$R_n^{(-1)}$: coefficients in the logarithm of a positive asymptotic
quantity need not be positive.

\subsection{The missing orders one third and one quarter}
\begin{example}[Order one third]\label[example]{ex:third}
Take $\theta=1$, $w_j=(j!)^{1/3}$, and arbitrary $C>0$.
Then $a_1=C$, $a_2=2^{1/3}C$, and $a_3=6^{1/3}C$.
\Cref{thm:main} gives exactly \eqref{eq:introthird}. Equivalently,
\begin{align}
 \sup_{f\in\mathcal B_{1/3,A,C}}|[z^n]f\inv|
 ={}&CA^{n-1}((n-1)!)^{1/3}\notag\\
 &\times\exp\left\{Cn^{2/3}
       +\left(2^{1/3}C+\frac23C^2\right)n^{1/3}\right\}\notag\\
 &\times\exp\left\{6^{1/3}C+\frac53\,2^{1/3}C^2
                        +\frac79C^3\right\}(1+o(1)).
 \label{eq:thirdextreme}
\end{align}
This is a sharp supremum, attained at every degree by the same series.
\end{example}

\begin{example}[Order one quarter]\label[example]{ex:quarter}
For $\theta=1$, $w_j=(j!)^{1/4}$, and $a_j=C(j!)^{1/4}$,
\begin{align}
 \log R_n^{(1)}={}&a_1n^{3/4}
 +\left(a_2+\frac58a_1^2\right)n^{1/2}\notag\\
 &+\left(a_3+\frac32a_1a_2+\frac{21}{32}a_1^3\right)n^{1/4}\notag\\
 &+a_4+\frac74a_1a_3+a_2^2
             +\frac{11}{4}a_1^2a_2+\frac{15}{16}a_1^4+o(1).
 \label{eq:quarter}
\end{align}
The lower-order exponent terms cannot be recovered from the leading
law $\log R_n\sim Cn^{3/4}$. Each is required for a relative-error
asymptotic.
\end{example}

\subsection{Recovery of the predecessor's regimes}
For $1/2<s<1$, only $p_1n^{1-s}$ survives. With $w_1=1$,
\[
 R_n^{(1)}=\exp(Cn^{1-s})(1+o(1)).
\]
At $s=1/2$, \eqref{eq:repoweights} gives
$a_2=C\sqrt{\nu+2}$, so
\begin{equation}
 R_n^{(1)}=
 \exp\left\{C\sqrt n+C\sqrt{\nu+2}+\frac34C^2\right\}(1+o(1)).
 \label{eq:half}
\end{equation}
These agree with the predecessor \cite{subexp}; they are consistency
checks, not newly claimed special cases. The new finite formula
continues below that boundary without introducing a new argument for
each interval $1/(r+1)<s\le1/r$.

For the alternative coefficient ball
$|f_n|\le B_0 A^{n-1}(n!)^s$, use $\nu=1$ in
\eqref{eq:repoweights} and set $C=2^sB_0$.
This conversion changes the constants, even though the $n$th-root
Gevrey type is unchanged.

\section{Critical windows and finite-head universality}\label{sec:universality}

\subsection{Every transition on its natural scale}
\begin{theorem}[All reciprocal critical windows]\label{thm:windows}
Fix a positive integer $r$, $\theta\ge-1$, and a family of admissible
weights and positive constants satisfying the uniform hypotheses of
\cref{sec:main} in a neighborhood of $s=1/r$. Write
$p_k(s)$ for the corresponding pressure coefficients, including the
$s$-dependence of the finite weights. For
\begin{equation}
 s_n=\frac1r+\frac{t}{\log n},
 \label{eq:window}
\end{equation}
uniformly for $t$ in every bounded real interval,
\begin{equation}
 \log R_n^{(\theta)}(s_n)
 -\sum_{k=1}^{r-1}p_k(s_n)n^{1-ks_n}
 \longrightarrow p_r(1/r)\ee^{-rt}.
 \label{eq:windowlimit}
\end{equation}
Thus the ratio obtained by removing the lower-order exponential terms
converges to $\exp(p_r(1/r)\ee^{-rt})$.
\end{theorem}
\begin{proof}
Choose a small fixed neighborhood of $1/r$ on which
$(r+1)s>1$, and use the uniform theorem with $J=r$.
After the indicated subtraction, the expression is
$p_r(s_n)n^{1-rs_n}+o(1)$. The power is exactly $\ee^{-rt}$,
and continuity gives $p_r(s_n)\to p_r(1/r)$ uniformly for bounded $t$.
\end{proof}

\begin{remark}[The lower-order coefficients must move]
In \eqref{eq:windowlimit}, replacing $p_k(s_n)$ by $p_k(1/r)$ for
$k<r$ is generally wrong. An $O(1/\log n)$ coefficient error multiplied
by $n^{1-k/r}$ need not vanish. The formula retains the actual moving
coefficients and therefore resolves the constant-scale transition
without hiding a larger error.
\end{remark}

At $r=1$ and $w_1=1$, \eqref{eq:windowlimit} becomes
$R_n\to\exp(C\ee^{-t})$. At $r=2$, after removing
$Cn^{1-s_n}$, the reversion limit is
$\exp((C\sqrt{\nu+2}+3C^2/4)\ee^{-2t})$.
For every larger $r$, \eqref{eq:criticalcoefficient} computes the
corresponding constant by a finite polynomial calculation.

\subsection{The precise number of weights that survive}
\begin{theorem}[Finite-head universality]\label{thm:universality}
Fix $s>0$, $C>0$, and $\theta\ge-1$. Let $w$ and $\widetilde w$ be
two admissible weight sequences with the same exponent $s$, but
possibly different tail scales, shifts, and finite heads. Set
$K=\lfloor1/s\rfloor$.
If
\begin{equation}
 w_j=\widetilde w_j\qquad(1\le j\le K),
 \label{eq:headagree}
\end{equation}
then their individually normalized inverse coefficients satisfy
\begin{equation}
 \frac{R_n^{(\theta)}(w)}{R_n^{(\theta)}(\widetilde w)}\longrightarrow1.
 \label{eq:universalratio}
\end{equation}
For $s>1$, the condition is empty and the conclusion still holds.
\end{theorem}
\begin{proof}
For $s\le1$, apply \cref{thm:main} with $J=K$. Every surviving
coefficient is a polynomial in the same initial weights, so the two
logarithms differ by $o(1)$. Exponentiation gives
\eqref{eq:universalratio}. For $s>1$, both normalized sequences tend
to $1$.
\end{proof}

The individual normalization is important. The theorem compares
$[z^n]G_\theta/(Cw_{n-1})$ with its counterpart for $\widetilde w$;
it does not claim equality of the unnormalized inverse coefficients.
Changing the factorial shift can produce a power of $n$ in the ratio
of the baseline weights.

\begin{proposition}[Sharp sensitivity to the first changed weight]
\label[proposition]{prop:sensitivity}
Suppose $w_j=\widetilde w_j$ for $j<r$ and
$w_r\ne\widetilde w_r$. Let $\Delta=C(w_r-\widetilde w_r)$.
If $rs<1$, then
\begin{equation}
 \log\frac{R_n^{(\theta)}(w)}{R_n^{(\theta)}(\widetilde w)}
 =\Delta n^{1-rs}+o(n^{1-rs}).
 \label{eq:sensitivity}
\end{equation}
If $rs=1$, that ratio tends to $\ee^\Delta$.
If $r>1/s$, it tends to $1$.
\end{proposition}
\begin{proof}
Formula \eqref{eq:pk} has the triangular form
\begin{equation}
 p_k=a_k+\text{a polynomial in }a_1,\ldots,a_{k-1}.
 \label{eq:triangular}
\end{equation}
Indeed, the $a_k t^k$ term can enter only through the linear term
of $\Psi$ and the leading denominator term in $B$, giving
$ka_k/k=a_k$. Thus the first pressure difference is $\Delta$ at
index $r$. Later powers of $n$ are strictly smaller. Apply the main
theorem and distinguish the three cases.
\end{proof}

\begin{example}[Tail changes invisible after normalization]
For a fixed $s=1/3$, keep $w_1,w_2,w_3$ fixed and replace all sufficiently
large weights by any $D_*\Gamma(j+\nu+1)^{1/3}$, with arbitrary
positive intervening finite values. The normalized reversion ratio to
the original sequence tends to $1$. Changing only $w_3$ by $\delta$
instead multiplies the limiting ratio by $\exp(C\delta)$.
This separates a genuinely observable head change from an invisible
tail change at the chosen normalization.
\end{example}

\section{Explicit Poisson centers for the small parts}\label{sec:distribution}

\subsection{The probability model attached to the finite sum}
Assign to each ordered composition $(j_1,\ldots,j_k)$ of $m=n-1$
probability proportional to its positive contribution in
\eqref{eq:exactR}. Normalization by $R_n^{(\theta)}$ makes this a
probability distribution. When there is a unique part at least $m-h$,
remove it and let $M_{j,n}$ count the remaining parts of size $j$.
On the exceptional event, define all these counts to be zero; the
choice on that event has no effect on the limits below.

At $\theta=-1$, grouping outdegrees in the Lagrange formula for
\eqref{eq:treeeq} gives exactly the same distribution of the nonzero
outdegree counts. Thus the statement can also be read as a theorem
about degree counts in a weighted plane tree. One can verify the identification by inserting an independent marking
variable for each outdegree into the root equation before applying
Lagrange inversion. No assertion about the location of the large vertex
is needed.

\begin{theorem}[An explicit product-Poisson approximation]
\label{thm:Poissonlaw}
Under the hypotheses of \cref{thm:main}, with $J$ fixed as there,
let $Y_{j,n}$ be independent Poisson variables with means
\begin{equation}
 \mu_{j,n}=n\frac{a_jy(n^{-s})^j}{1-\theta A(y(n^{-s}))},
 \qquad1\le j\le J.
 \label{eq:explicitcenters}
\end{equation}
Then
\begin{equation}
 \TV\left(\Law((M_{j,n})_{j=1}^J),
                \bigotimes_{j=1}^J\Pois(\mu_{j,n})\right)\to0.
 \label{eq:TV}
\end{equation}
With probability tending to one, every part other than the unique
large part has size at most $\lfloor1/s\rfloor$. If $s>1$, there are
no other parts with probability tending to one.
\end{theorem}
\begin{proof}
\Cref{lem:localization,lem:finitecut} show that the exceptional
configurations and those containing a small part of size greater than
$J$ have total relative mass $o(1)$.
On the remaining multiplicity vectors, \cref{lem:product} changes the
unnormalized mass by a uniform factor $1+o(1)$.
The exact Poisson change of measure in \eqref{eq:tiltexact} gives the
Radon--Nikodym factor proportional to
$\exp(\mathcal E_n)\ind_{D\le h}$ relative to the product-Poisson law.
By \eqref{eq:L1tilt}, this factor tends to $1$ in $L^1$, and its
normalizing constant tends to $1$. This proves \eqref{eq:TV}.

If $js>1$, then $\mu_{j,n}\sim a_jn^{1-js}\to0$.
Together with the finite-size cutoff this removes all sizes greater
than $\lfloor1/s\rfloor$. When $s>1$ even the first mean tends to
zero, so the cloud is empty with probability tending to one.
\end{proof}

\subsection{Gaussian and critical Poisson limits}
For $js<1$, the means diverge, and the ordinary Poisson central limit
theorem plus \eqref{eq:TV} gives
\begin{equation}
 \frac{M_{j,n}-\mu_{j,n}}{\sqrt{\mu_{j,n}}}
       \ \xrightarrow{\ d\ }\ \mathcal N(0,1).
 \label{eq:CLT}
\end{equation}
For $js=1$,
\begin{equation}
 M_{j,n}\ \xrightarrow{\ d\ }\ \Pois(a_j).
 \label{eq:criticalPoisson}
\end{equation}
All these limits hold jointly with independent components. These are
statements about distributions. Total-variation convergence by itself
does not assert convergence of unbounded moments; we call
$\mu_{j,n}$ \emph{centering parameters}, without silently identifying
them with the exact expectations of $M_{j,n}$.

\subsection{All centering corrections by finite arithmetic}
Expanding the small solution of \eqref{eq:y} gives
\begin{equation}
 y=x+sa_1x^2+O(x^3).
\end{equation}
Consequently
\begin{equation}
 \mu_{j,n}=a_jn^{1-js}
 \left(1+(\theta+js)a_1n^{-s}+O(n^{-2s})\right).
 \label{eq:firstcenter}
\end{equation}
For small $s$, the correction in this formula may itself exceed the
Gaussian standard deviation, so leading-order centering can fail.
The finite implicit equation supplies every further coefficient.
More explicitly, if
\begin{equation}
 \beta_j(x)=\sum_{k\ge j}b_{j,k}x^k,
\end{equation}
then Lagrange inversion gives
\begin{equation}
 b_{j,k}=\frac1k[t^{k-1}]
 \frac{d}{dt}\left(\frac{a_jt^j}{1-\theta A(t)}\right)
            (1-B(t))^{-sk}.
 \label{eq:centeringcoefficients}
\end{equation}
This is again a finite coefficient calculation for every selected $k$.

\begin{corollary}[How many centering terms suffice]\label[corollary]{cor:centering}
Suppose $js<1$. If an integer $M\ge j$ satisfies
\begin{equation}
 (2M+2-j)s>1,
 \label{eq:centerdepth}
\end{equation}
then replacing $\mu_{j,n}$ in the numerator of \eqref{eq:CLT} by
\begin{equation}
 \widehat\mu_{j,n}=\sum_{k=j}^M b_{j,k}n^{1-ks}
 \label{eq:finitecenter}
\end{equation}
does not change the limiting standard normal law.
\end{corollary}
\begin{proof}
The omitted tail is $O(n^{1-(M+1)s})$.
The standard deviation is asymptotic to
$\sqrt{a_j}\,n^{(1-js)/2}$. Condition \eqref{eq:centerdepth} says
exactly that the former order is smaller than the latter.
\end{proof}

The $j=1$, $M=1$ condition is $3s>1$; the next correction permits
$5s>1$. Thus the hierarchy of increasingly necessary centering terms
has its own thresholds. Formula \eqref{eq:centeringcoefficients}
replaces separate ad hoc calculations at each such threshold.

\section{The weighted-tree specialization}\label{sec:trees}

Give every plane rooted tree $T$ weight
\begin{equation}
 \operatorname{wt}(T)=\prod_{v\in T}q_{\deg^+(v)},
 \qquad q_0=1,\quad q_j=Cw_j\ (j\ge1),
\end{equation}
and let $Z_n$ be the sum of weights over trees with $n$ vertices.
The standard root decomposition proves formally that
$T(z)=\sum_{n\ge1}Z_nz^n$ satisfies \eqref{eq:treeeq}.
Hence $Z_n=[z^n]G_{-1}$, and the main theorem proves
\begin{equation}
 Z_n=Cw_{n-1}\exp\left\{\sum_{k=1}^J
                 p_k(s,-1;\mathbf a)n^{1-ks}\right\}(1+o(1)).
 \label{eq:treeasymptotic}
\end{equation}

With $C=1$ and $w_j=(j!)^s$, the first two exponent terms are
\begin{equation}
 n^{1-s}+\left(2^s-\frac{1-s}{2}\right)n^{1-2s}.
 \label{eq:treefirsttwo}
\end{equation}
This agrees with the published correction in \cite[Remark~2.8]{jjs}.
The explicit formula \eqref{eq:pk} continues it to every required order.
For instance, at $s=1/3$ it yields
\begin{align}
 \log\frac{Z_n}{((n-1)!)^{1/3}}
 ={}&n^{2/3}+\left(2^{1/3}-\frac13\right)n^{1/3}\notag\\
 &+6^{1/3}-\frac13\,2^{1/3}+\frac19+o(1).
 \label{eq:treethird}
\end{align}
The corresponding extremal reversion constant in
\eqref{eq:introthird} has different interaction terms. The
interpolation parameter explains both without equating them.

Likewise, \eqref{eq:firstcenter} becomes
\[
 \mu_{j,n}=j!^s n^{1-js}
       \left(1-(1-js)n^{-s}+O(n^{-2s})\right),
\]
consistent with the known first centering correction
\cite[Theorem~2.6]{jjs}. The new coefficient formula
\eqref{eq:centeringcoefficients} specifies as many further terms as
\eqref{eq:centerdepth} requires. The tree conclusions concern degree
counts and normalization; additional geometric information about branch
shapes or the location of the large vertex is not asserted here.

\section{Computation, reproducibility, and interpretation}\label{sec:computation}

\subsection{A coefficient recurrence without enumerating compositions}
For a fixed target degree $n$, define
\begin{equation}
 H_n(z)=(1-\theta CW(z))^{-n/\theta}
        =\sum_{k\ge0}B_{n,k}z^k,
 \label{eq:Hn}
\end{equation}
with the exponential interpretation at $\theta=0$.
Formal differentiation gives
$(1-\theta CW)H_n'=nCW'H_n$. Therefore
\begin{equation}
 B_{n,0}=1,\qquad
 kB_{n,k}=C\sum_{j=1}^k
       \bigl(nj+\theta(k-j)\bigr)w_jB_{n,k-j},
 \label{eq:recurrence}
\end{equation}
and
\begin{equation}
 R_n^{(\theta)}=\frac{B_{n,n-1}}{nCw_{n-1}}.
 \label{eq:recurrenceR}
\end{equation}
For $k\le n-1$ and $\theta\ge-1$, every multiplier in
\eqref{eq:recurrence} is positive. A single coefficient can be
computed in $O(n^2)$ arithmetic operations and $O(n)$ stored numbers.
This is an arithmetic-operation count, not a bit-complexity estimate:
exact integer and rational entries can become large.

For rational $s$, $\theta$, and finite $a_j$, formula \eqref{eq:pk}
produces exact rational pressure coefficients whenever the $a_j$
are rational. Factorial weights at nonintegral $s$ need not themselves
be rational; the exact tests deliberately use rational finite inputs
to check the universal algebra, and do not relabel floating-point
factorial powers as exact values.

\subsection{What the supplied checks actually verify}
The standard-library script \path{code/verify.py} records 1,066 scalar
checks. It verifies both inverse compositions through degree $18$ in
12 parameter cases, compares ordered-composition sums with grouped
cloud sums, and checks the pressure formula independently through
degree $7$ in 20 rational cases. It also tests the displayed first four
polynomials, the critical pure-$a_1$ formula in 50 cases through
$r=10$, and all 256 sign patterns in a degree-nine extremal box.

The second script, \path{code/diagnostics.py}, computes normalized
coefficients by a log-scaled version of \eqref{eq:recurrence}. Its
45 recorded rows use $s\in\{1/2,2/5,1/3,1/4,1/5\}$,
$\theta\in\{-1,0,1\}$, and $n\in\{100,1000,10000\}$, with
$C=1$, $\nu=0$. Fifteen degree-$70$ values are independently recomputed
at 80 decimal digits. The largest absolute discrepancy between those
logarithms and the ordinary floating-point computation was below
$6.4\times10^{-14}$.

These comparisons test implementations and finite identities. They
are not interval certificates, a proof of the limits, or evidence that
an asymptotic regime has started at a particular finite degree. The
analytic proofs, rather than the numerical fit, support the theorems.

\subsection{Visible preasymptotic behavior}
Let $S_n=\sum_{k=1}^{\lfloor1/s\rfloor}p_kn^{1-ks}$ for the
reversion endpoint. The following table uses $C=1$, $\nu=0$.
The last column uses the analytic pressure from the same finite
polynomial $A$. A dash means that the diagnostic routine did not locate
the small positive real branch at that finite argument; the theorem
only guarantees that branch for all sufficiently large $n$.

\begin{table}[htbp]
\centering\small
\begin{tabular}{rrrrr}
\toprule
$s$ & $n$ & $\log R_n$ & $\log R_n-S_n$
 & $\log R_n-nP(n^{-s})$\\
\midrule
$1/3$ & 100 & 56.504158 & 21.322616 & \textemdash\\
$1/3$ & 1,000 & 126.387080 & 2.426436 & 0.721743\\
$1/3$ & 10,000 & 511.209609 & 0.848884 & 0.239161\\
$1/4$ & 100 & 83.728129 & 10.723605 & \textemdash\\
$1/4$ & 1,000 & 342.874419 & 74.579330 & 56.517991\\
$1/4$ & 10,000 & 1237.289700 & 5.243241 & 0.738160\\
\bottomrule
\end{tabular}
\caption{Non-certified logarithmic diagnostics. A residual need not be
monotone and need not be small at the displayed degrees.}
\label{tab:diagnostics}
\end{table}

The order-one-quarter residual grows sharply between $n=100$ and
$n=1000$ before decreasing. In particular, it would be misleading to
infer the constant term from a low-degree numerical fit. Even at
$n=10000$, a logarithmic error $5.24$ corresponds to a large relative
coefficient error. The resummed pressure often reduces that error in
these examples, but no universal finite-degree superiority theorem is
claimed.

\subsection{Reproduction and formalization boundary}
From the package directory, the exact checks require only Python's
standard library:
\begin{verbatim}
python code/verify.py
\end{verbatim}
The optional diagnostic run additionally requires NumPy, SciPy, and
mpmath:
\begin{verbatim}
python code/diagnostics.py --max-n 10000
\end{verbatim}
The default scripts reproduce the recorded output paths; use their
\texttt{--output} option for a separate destination. For the exact
checker this option is a JSON file path; for diagnostics it is a
directory. The manuscript is standalone and does not require generated
table files to compile:
\begin{verbatim}
pdflatex -interaction=nonstopmode -halt-on-error article.tex
pdflatex -interaction=nonstopmode -halt-on-error article.tex
pdflatex -interaction=nonstopmode -halt-on-error article.tex
\end{verbatim}

A future proof-assistant implementation has a useful division of labor.
The exact Lagrange sum, grouping identity, pressure-coefficient formula,
and finite coefficient computations form an algebraic layer. The
localization bounds, cutoff insertion lemma, Poisson tail control, and
uniform asymptotic statements form a separate analytic layer. The
present package supplies no Lean source and makes no kernel-verification
claim about either layer.

\section{Further research questions}\label{sec:future}

\subsection{Effective relative-error bounds}
The proofs establish $o(1)$, not a sharp computable onset. Derive an
explicit error bound for
$R_n^{(\theta)}\exp(-nP(n^{-s}))-1$ in terms of $s$, the finite head,
and the tail parameters. A useful answer would separate the discarded
size tail, the Euler-sum correction, and the quadratic Poisson tilt.
The diagnostics show why an effective result would matter in practice.

\subsection{The first vanishing corrections}
The present pressure gives every nonvanishing term in $\log R_n$.
Determine the next terms, including the shift $\nu$, the one-step
Euler corrections, and the determinant or cumulant corrections from
the Poisson tilt. The natural scales include both powers of $n^{-s}$
and ordinary inverse powers of $n$, with coincidences at rational $s$.
A uniform calculus should handle these resonances without duplicating
terms.

\subsection{The joint small-order limit}
Here $s$ is fixed away from zero in every uniform statement.
When $s=s_n\downarrow0$, the required number of small sizes diverges.
Identify the crossover between finite-dimensional condensation and a
regime with an unbounded collection of active sizes. Simply taking
$J=\lfloor1/s_n\rfloor$ in the present proof is not justified because
its constants and analytic neighborhoods depend on $J$.

\subsection{Beyond an exact factorial tail}
Replace \eqref{eq:tail} by a local-quotient hypothesis such as
$w_{n-1}/w_n=n^{-s}(1+\varepsilon_n)$. What quantitative control of
$\varepsilon_n$ over windows of length $O(n^{1-s})$ preserves the same
pressure? A pointwise ratio limit alone may hide accumulated corrections.
Regularly varying modifications and oscillatory perturbations should
be distinguished rather than treated by a single informal heuristic.

\subsection{Vanishing head weights and lattice restrictions}
Allow some $w_j$ to be zero, with smallest positive index $r_0>1$.
The first cloud scale then changes from $n^{1-s}$ to
$n^{1-r_0s}$. Determine the exact analogue of the pressure theorem,
including congruence restrictions and possible zero inverse
coefficients. This would extend the present strictly positive class
without assuming that missing sizes can be ignored harmlessly.

\subsection{Signs and cancellation beyond the extremum}
For a general signed or complex coefficient series, the positive
Lagrange model no longer represents the inverse coefficient in absolute
value. Find structural noncancellation conditions under which a
finite-jet equivalent remains valid. Conversely, classify finite-head
changes that can suppress an entire leading exponential scale.
The coefficient-ball supremum proved here does not answer this question.

\subsection{Several variables and operator coefficients}
Extend the fixed-radius, fixed-prefactor extremal problem to formal
maps of several variables or to operator-valued coefficients.
Type-level stability is weaker than the multiplicative control sought
here. A useful theorem would identify which finite tensors replace the
scalar head $a_1,\ldots,a_K$ and whether a finite-dimensional pressure
still captures all surviving subexponential terms.

\subsection{The positivity boundary in the interpolation parameter}
For $\theta<-1$, the factors in \eqref{eq:exactR} can change sign.
Determine whether a modified theorem survives in a restricted parameter
range, whether new cancellation transitions arise, and which algebraic
parts of \eqref{eq:pk} retain an asymptotic interpretation. Formal
coefficient identities alone do not settle that analytic issue.

\subsection{Quantitative probability and geometric information}
Obtain rates in \eqref{eq:TV}, moderate deviations for the total cloud
size, and sufficiently strong tail control to compare actual moments
with the explicit centering parameters. At the tree endpoint, connect
the pressure derivatives to geometric observables such as branch
profiles. These questions go beyond the count distribution proved here.

\subsection{Certified computation and proof-assistant integration}
Develop a certified interval solver for the small branch of
\eqref{eq:y}, with a proof that the chosen branch is the analytic germ
required by the theorem. Formalize the finite-jet extraction and the
cutoff Poisson lemma, then connect them to an independently verified
Lagrange inversion library. A complete implementation should report
separately algebraic exactness, analytic error bounds, and machine
verification.

\section{Conclusion}
A fixed Gevrey coefficient ball contains more asymptotic information
than its coefficient type. The missing information is finite but grows
as the order decreases: at order $s$, the initial
$\lfloor1/s\rfloor$ weights and their interactions determine every
nonvanishing exponent term. The formula
\[
 p_k=\frac1k[t^k]\Psi_{ks}\!\left(\frac{tA'(t)}{1-\theta A(t)}\right)
\]
computes those terms directly, and the finite analytic pressure
packages them into a single local equation.

This supplies a full multiplicative answer to the predecessor's
unresolved $0<s<1/2$ extremal-reversion problem, while also giving
explicit critical windows and centering formulas in the weighted-tree
specialization. The decisive steps are the relative finite-size
cutoff and the cutoff-preserving Poisson tilt. Neither requires the
divergent generating function to acquire an analytic meaning.

\begin{thebibliography}{9}
\small\setlength{\itemsep}{3pt}\setlength{\parskip}{0pt}
\bibitem{repo}
Vladimir Reshetnikov,
\emph{ProveIt}, mathematical formalizations and research corpus.
Snapshot \snapshot, inspected 30 September 2026.
\href{https://github.com/VladimirReshetnikov/ProveIt/tree/a4268e78ebd0f6bf71ba609c07f4b3415748f532}{Pinned repository snapshot}.

\bibitem{subexp}
\emph{Sharp Subexponential Reversion: Exact extremal cost, a Gevrey
transition, and composition of two divergent series}.
Research article prepared for Vladimir Reshetnikov, 29 September 2026,
with repository editorial amendments.
ProveIt package
\path{Sharp_Subexponential_Cost_Gevrey_Reversion/}, under
\path{Analysis/Transseries/docs/series-and-transseries/}.
\href{https://github.com/VladimirReshetnikov/ProveIt/blob/a4268e78ebd0f6bf71ba609c07f4b3415748f532/Analysis/Transseries/docs/series-and-transseries/Sharp_Subexponential_Cost_Gevrey_Reversion/article.tex}{Pinned article source}.
The README and the extremal, localization, and multiplicative-asymptotic
sections were consulted.

\bibitem{weighted}
\emph{Sharp Weighted Type Under Formal Reversion: Zero-loss inversion,
exact Borel radii, and non-D-finiteness in countable exponential feedback}.
Research article prepared for Vladimir Reshetnikov, 29 September 2026,
with repository editorial amendments.
ProveIt package \path{Sharp_Weighted_Type_Formal_Reversion/},
under the same series-and-transseries directory.
\href{https://github.com/VladimirReshetnikov/ProveIt/blob/a4268e78ebd0f6bf71ba609c07f4b3415748f532/Analysis/Transseries/docs/series-and-transseries/Sharp_Weighted_Type_Formal_Reversion/README.md}{Pinned package description}.

\bibitem{gessel}
Ira M. Gessel,
\emph{Lagrange Inversion},
arXiv:1609.05988, 2016.
\href{https://arxiv.org/abs/1609.05988}{Article and formal coefficient formulas}.

\bibitem{jjs}
Svante Janson, Thordur Jonsson, and Sigurdur \"Orn Stef\'ansson,
\emph{Random trees with superexponential branching weights},
Journal of Physics A: Mathematical and Theoretical
\textbf{44} (2011), 485002.
\href{https://doi.org/10.1088/1751-8113/44/48/485002}{doi:10.1088/1751-8113/44/48/485002}.
\href{https://arxiv.org/abs/1104.2810}{arXiv:1104.2810};
Theorems 2.5--2.6 and Remarks 2.7--2.8 in the inspected version.

\bibitem{borinsky}
Michael Borinsky,
\emph{Generating asymptotics for factorially divergent sequences},
Electronic Journal of Combinatorics \textbf{25}(4) (2018), Paper P4.1.
\href{https://doi.org/10.37236/5999}{doi:10.37236/5999}.
\href{https://arxiv.org/abs/1603.01236}{arXiv:1603.01236}.
\end{thebibliography}
\end{document}
